\section{Datenschutzrecht (DSG)}

\subsection{Grundbegriffe und Geltungsbereich}

\begin{definition}
    \textbf{Geltungsbereich des Schweizer Datenschutzgesetzes (revDSG):} Das revidierte DSG bezweckt den Schutz der Persönlichkeit und der Grundrechte natürlicher Personen, über die Personendaten bearbeitet werden (Art. 1 DSG). Juristische Personen fallen seit der Revision 2023 nicht mehr unter den Schutzbereich des DSG.
\end{definition}

\begin{definition}
    \textbf{Kategorien von Daten:}
    \begin{itemize}
        \item \textbf{Personendaten (Art. 5 lit. a DSG):} Alle Angaben, die sich auf eine bestimmte oder bestimmbare natürliche Person beziehen. Bestimmbar ist eine Person, wenn sie direkt oder über Zusatzwissen (z.B. Verknüpfung mit Logfiles, Cookie-IDs, dynamischen IP-Adressen oder Verhaltensmustern) identifiziert werden kann.
        \item \textbf{Besonders schützenswerte Personendaten (Art. 5 lit. c DSG):} Daten über:
        \begin{itemize}
            \item religiöse, weltanschauliche, politische oder gewerkschaftliche Ansichten oder Tätigkeiten,
            \item die Gesundheit, die Intimsphäre oder die Zugehörigkeit zu einer Rasse oder Ethnie,
            \item genetische Daten und biometrische Daten, die eine Person eindeutig identifizieren,
            \item verwaltungs- und strafrechtliche Verfolgungen oder Sanktionen,
            \item Massnahmen der sozialen Hilfe.
        \end{itemize}
    \end{itemize}
\end{definition}

\begin{remark}
    \textbf{Anonymisierung vs. Pseudonymisierung:}
    \begin{itemize}
        \item \textbf{Anonymisierung:} Der Personenbezug wird irreversibel vernichtet, sodass die betroffene Person von niemandem mehr identifiziert werden kann. Vollständig anonymisierte Daten fallen nicht mehr unter das DSG.
        \item \textbf{Pseudonymisierung:} Der Name oder Identifikator wird durch ein Pseudonym ersetzt (z.B. User-ID), wobei der Zuordnungsschlüssel getrennt aufbewahrt wird. Da der Personenbezug mit dem Schlüssel wiederhergestellt werden kann, bleiben pseudonymisierte Daten vollwertige Personendaten im Sinne des DSG.
        \item \textbf{Re-Identifikationsrisiko bei Big Data:} Bei umfangreichen Datenmengen genügt oft die Korrelation weniger Merkmale (z.B. PLZ, Geburtsdatum, Geschlecht), um scheinbar anonyme Datensätze eindeutig konkreten Personen zuzuordnen.
    \end{itemize}
\end{remark}

\subsection{Bearbeitungsgrundsätze (Art. 6 und 7 DSG)}

\begin{theorem}
    \textbf{Zwingende Bearbeitungsgrundsätze (Art. 6 DSG):} Jede Datenbearbeitung muss den folgenden Grundsätzen genügen:
    \begin{enumerate}
        \item \textbf{Rechtmässigkeit:} Datenbearbeitungen dürfen nicht gegen Rechtsnormen verstossen.
        \item \textbf{Treu und Glauben (Transparenz):} Die Datenbeschaffung und deren Zweck müssen für die betroffene Person erkennbar sein; keine heimliche oder irreführende Datenbeschaffung.
        \item \textbf{Zweckbindung:} Personendaten dürfen nur zu dem Zweck beschafft werden, der bei der Beschaffung angegeben wurde, aus den Umständen ersichtlich oder gesetzlich vorgesehen ist.
        \item \textbf{Verhältnismässigkeit:} Die Bearbeitung muss zur Erreichung des Zwecks geeignet und erforderlich sein (Datenminimierung: nur so viele Daten wie nötig bearbeiten und so früh wie möglich löschen).
        \item \textbf{Richtigkeit:} Wer Personendaten bearbeitet, muss sich über deren Richtigkeit vergewissern; unrichtige Daten müssen korrigiert oder gelöscht werden.
        \item \textbf{Datensicherheit (Art. 8 DSG):} Angemessene technische und organisatorische Massnahmen (TOMs) gegen unbefugte Bearbeitung, Verlust oder Zerstörung.
    \end{enumerate}
\end{theorem}

\begin{important}
    \textbf{Legalitätsprinzip und Opt-out im Schweizer Recht:} \\
    Die Bearbeitung von Personendaten ist grundsätzlich \textbf{erlaubt}, solange kein Bearbeitungsgrundsatz verletzt wird und kein Widerspruch der betroffenen Person vorliegt.
    
    Eine explizite Einwilligung (Opt-in) ist nach Schweizer Recht nur erforderlich bei besonders schützenswerten Daten, hochriskantem Profiling oder wenn eine Persönlichkeitsverletzung vorliegt, die nicht durch ein überwiegendes privates/öffentliches Interesse oder Gesetz gerechtfertigt ist.
\end{important}

\subsection{Rollenverteilung und Datenübermittlung}

\begin{definition}
    \textbf{Verantwortlicher vs. Auftragsbearbeiter:}
    \begin{itemize}
        \item \textbf{Verantwortlicher (Controller, Art. 5 lit. j DSG):} Die private Person oder das Bundesorgan, die allein oder zusammen mit anderen über den Zweck und die Mittel der Bearbeitung entscheidet. Trägt die Hauptverantwortung für die Einhaltung des DSG.
        \item \textbf{Auftragsbearbeiter (Processor, Art. 5 lit. k DSG):} Bearbeitet Personendaten ausschliesslich im Auftrag und nach Weisung des Verantwortlichen (z.B. Cloud-Provider, SaaS-Anbieter).
        \item \textbf{Auftragsbearbeitungsvertrag (AVV / DPA, Art. 9 DSG):} Die Delegation erfordert eine vertragliche oder gesetzliche Vereinbarung, die sicherstellt, dass der Auftragsbearbeiter Daten nur so bearbeitet, wie es der Verantwortliche selbst tun dürfte, und die Datensicherheit gewährleistet. Sub-Auftragsbearbeiter bedürfen der vorherigen Genehmigung.
    \end{itemize}
\end{definition}

\begin{theorem}
    \textbf{Bekanntgabe von Personendaten ins Ausland (Art. 16 f. DSG):}
    Personendaten dürfen grundsätzlich nur dann ins Ausland bekanntgegeben werden, wenn der Zielstaat einen \textbf{angemessenen Datenschutz} gewährleistet (Staatenliste des Bundesrats, z.B. alle EU/EWR-Staaten sowie das Swiss-US Data Privacy Framework).
    
    Fehlt ein Angemessenheitsbeschluss (z.B. USA ohne DPF, sonstige Drittstaaten), ist ein Export nur zulässig bei:
    \begin{itemize}
        \item Geeigneten Garantien, insbesondere durch vom EDÖB anerkannte \textbf{Standardvertragsklauseln (EU Standard Contractual Clauses mit Schweizer Anpassung)}, oder
        \item Ausnahmetatbeständen im Einzelfall (ausdrückliche Einwilligung, unmittelbare Vertragserfüllung).
    \end{itemize}
\end{theorem}

\subsection{Betroffenenrechte und Organisatorisches}

\begin{definition}
    \textbf{Rechte der betroffenen Personen:}
    \begin{itemize}
        \item \textbf{Auskunftsrecht (Art. 25 DSG):} Jede Person kann unentgeltlich verlangen, ob und welche Personendaten über sie bearbeitet werden, zu welchem Zweck und an wen sie übermittelt wurden (Frist: in der Regel innert 30 Tagen).
        \item \textbf{Berichtigung und Löschung:} Recht auf Berichtigung unrichtiger Daten oder Löschung unrechtmässig bearbeiteter Daten.
        \item \textbf{Recht auf Datenherausgabe und -übertragung (Datenportabilität, Art. 28 DSG):} Recht auf Aushändigung der eigenen Daten in einem gängigen elektronischen Format oder Übertragung an einen Drittanbieter.
    \end{itemize}
\end{definition}

\begin{definition}
    \textbf{Organisatorische Unternehmenspflichten:}
    \begin{itemize}
        \item \textbf{Datenschutzerklärung (DSE, Art. 19 DSG):} Umfassende Informationspflicht bei der Beschaffung von Personendaten (Identität, Zweck, Empfänger, Auslandstransfers).
        \item \textbf{Datenschutz-Folgenabschätzung (DSFA, Art. 22 DSG):} Pflicht zur vorgängigen schriftlichen Risikoanalyse bei Bearbeitungen, die ein hohes Risiko für die Persönlichkeit mit sich bringen (z.B. flächendeckende Videoüberwachung, Scoring, KI-Profiling).
        \item \textbf{Privacy by Design \& Privacy by Default (Art. 7 DSG):} Systeme müssen ab Entwicklung datenschutzfreundlich konzipiert sein; Voreinstellungen müssen auf das Minimum an Datenbearbeitung beschränkt sein.
        \item \textbf{Meldepflicht bei Datenschutzverletzungen (Data Breaches, Art. 24 DSG):} Verletzungen der Datensicherheit mit hohem Risiko für die Betroffenen müssen unverzüglich dem Eidgenössischen Datenschutz- und Öffentlichkeitsbeauftragten (EDÖB) gemeldet werden.
    \end{itemize}
\end{definition}

\begin{important}
    \textbf{Strafbestimmungen (Art. 60 ff. DSG):}
    Bei vorsätzlichen Verstössen gegen zentrale DSG-Vorschriften drohen \textbf{persönliche Bussen bis zu CHF 250'000.--}:
    \begin{itemize}
        \item Bestraft wird die handelnde \textbf{natürliche Person} (z.B. Geschäftsführer, IT-Leiter, Datenschutzbeauftragter), nicht das Unternehmen!
        \item Strafbar sind u.a.: Falsche oder verweigerte Auskunft (Art. 25 DSG), Verletzung der Informationspflicht (Art. 19 DSG), Missachtung von Sorgfaltspflichten bei der Auftragsbearbeitung oder beim Auslandstransfer sowie Verletzung der beruflichen Schweigepflicht.
    \end{itemize}
\end{important}

\subsection{Sonderfälle und Begleitgesetze}

\begin{definition}
    \textbf{Spam-Verbot im E-Commerce (Art. 3 Abs. 1 lit. o UWG):}
    Wer Massenwerbung per E-Mail, SMS oder Messenger versendet, handelt unlauter, wenn nicht vorgängig die \textbf{Einwilligung der Kunden} eingeholt wurde (\textbf{Opt-in-Pflicht}), der korrekte Absender nicht genannt wird oder keine einfache, kostenlose Ablehnungsmöglichkeit vorhanden ist.
    
    \textbf{Ausnahme (Kundenbeziehung):} Keine vorherige Einwilligung nötig, wenn:
    \begin{enumerate}
        \item der Händler die Kontaktadresse beim Verkauf von Waren oder Dienstleistungen erhalten hat,
        \item die Werbung ausschliesslich eigene, ähnliche Waren oder Dienstleistungen betrifft, und
        \item dem Kunden bereits bei der Erhebung und bei jeder Sendung eine klare Abmeldemöglichkeit geboten wird.
    \end{enumerate}
\end{definition}

\begin{remark}
    \textbf{Sterneintrag und Telefonwerbung (Art. 3 Abs. 1 lit. u UWG):}
    Unlauter handelt, wer den Vermerk im Telefonbuch (,,Sterneintrag``), dass ein Kunde keine Werbeanrufe wünscht, missachtet und Werbeanrufe tätigt.
\end{remark}

\begin{definition}
    \textbf{EU AI Act (Überblick und Relevanz für die Schweiz):}
    Die Europäische Union reguliert Künstliche Intelligenz nach einem risikobasierten Ansatz:
    \begin{itemize}
        \item \textbf{Inakzeptables Risiko:} Vollständig verboten (z.B. Social Scoring, manipulative Verhaltensbeeinflussung, biometrische Echtzeitüberwachung).
        \item \textbf{Hohes Risiko (High-Risk AI):} Strenge Zulassungsvoraussetzungen (Konformitätsbewertung, Risikomanagementsystem, hohe Datenqualität, Transparenz, menschliche Aufsicht; z.B. KI in kritischer Infrastruktur, HR-Bewerberscreening, Kreditwürdigkeitsprüfung).
        \item \textbf{Minimales/Geringes Risiko:} Transparenzpflichten (z.B. Kennzeichnung von KI-generierten Inhalten / Deepfakes, Offenlegung von Chatbots).
    \end{itemize}
    \textbf{Auswirkung auf Schweizer Unternehmen:} Schweizer Firmen unterliegen dem EU AI Act direkt, sobald sie KI-Systeme im EU-Markt in Verkehr bringen oder die Outputs ihrer KI in der EU genutzt werden (Marktortprinzip).
\end{definition}
